% ============================== ЗАКЛЮЧЕНИЕ ==============================
\chapter*{Заключение}
\addcontentsline{toc}{chapter}{Заключение}
\markboth{Заключение}{}

В изследването са изградени, обучени и експериментално оценени девет модела за откриване на прониквания в мрежовия трафик на център за данни: Random Forest, SVM с RBF ядро и XGBoost за DDoS, за Brute Force срещу FTP и SSH (BruteForce) и за Botnet (Bot). Всички етапи използват само реални записи от CSE-CIC-IDS2018, уловени в облачна инфраструктура на AWS: \val{rows-read} изтеглени и \val{rows-clean} след почистване. Векторът на признаците е редуциран в два етапа от \val{n-features} до \val{n-selected} признака. Моделите са оценени по TP, TN, FP, FN, Accuracy, Precision, Recall (Sensitivity), F1 и Specificity с тълкуване „на 100 случая“, допълнени с ROC-AUC, доверителни интервали, тест на McNemar и явна проверка за преобучение.

Всички девет модела постигат Accuracy над 99,5\%. При DDoS и BruteForce Random Forest и XGBoost са практически безгрешни. При Bot, който е най-трудната от трите атаки, XGBoost постига \val{XGB-Bot-Accuracy}, Random Forest \val{RF-Bot-Accuracy}, а SVM \val{SVM-Bot-Accuracy}. Разликите между точността при обучение, при кръстосана валидация и при тест са под \val{gap-max} пункта, т.е. моделите не са преобучени. За това допринасят регуляризацията на алгоритмите и правилото „една стандартна грешка“ при избора на хиперпараметри.

Предварителното изследване върху IoT набора CICIoT2023 показа, че Port Scan там не може да се открие с точност над около 85\%, независимо от алгоритъма и от методите за подобряване. Причината е, че признаците на отделния поток не съдържат отпечатъка на Port Scan. Сравнението на двата набора е основният извод: качеството на ML-базираната IDS се определя преди всичко от това дали представянето на трафика носи информация за атаката, а едва след това от избора на алгоритъм.

Сред алгоритмите XGBoost е най-подходящ за практическа IDS: той е най-точен, класифицира над милион записа в секунда при DDoS и дава обяснима важност на признаците. Random Forest е почти равностоен. SVM изостава по точност при Bot и по скорост и изисква подизвадка при обучение. Подборът на признаци по важност е по-подходящ от PCA: запазва същото качество с по-малко измерения и с интерпретируеми признаци.

Почти перфектните резултати са проверени с пет допълнителни анализа за приложимост. Без признаците, които отразяват TCP стека на атакуващия инструмент, точността почти не се променя. При хронологично разделяне DDoS и BruteForce остават над \val{rob-chrono-min-DDoSBF}. При Bot ансамблите от дървета пропускат около 10\% от по-късните атаки, а комбинираният детектор XGBoost или SVM ги открива почти всички (\val{rob-chronological-all-XGBorSVM-Bot-sens}). Анализите показват и две условия за практическото използване: обучаващите данни трябва да съдържат всеки инструмент на атаката и разнообразен нормален трафик от различни периоди, иначе фалшивите аларми върху непознат нормален трафик могат да достигнат десетки проценти.

Програмната реализация е декларативна и напълно възпроизводима. Атаките и алгоритмите са записи в регистри, всички числа в доклада се генерират от резултатите, а коректността на конвейера се проверява с автоматични тестове върху реални данни. Затова изследването може да се разшири с нови атаки, алгоритми, признаци и набори от данни с минимални промени.
